\documentclass[10pt]{article}
\usepackage[T1]{fontenc}
\usepackage[utf8]{inputenc}
\usepackage[french]{babel}
\usepackage{amsmath,amssymb,mathrsfs,mathpazo}
\usepackage{geometry}
\geometry{a4paper,margin=2cm}
\usepackage{graphicx,array,tabularx,tabularray,colortbl,multirow,multicol,booktabs,enumitem,pifont,ulem,fancybox,diagbox,cancel}
\usepackage{pstricks,pstricks-add,pst-plot,pst-tree,pst-node,pst-text,pst-func,pst-math,pst-bspline,pst-3dplot,pst-eucl,pst-solides3d}
\let\pscalculate\undefined
\usepackage{pst-bezier}
\usepackage{tikz,pgfplots,tkz-tab}
\usetikzlibrary{arrows,arrows.meta,patterns,calc,positioning,trees,decorations.markings,shapes.geometric,shapes.misc,shapes.arrows,intersections,decorations.pathmorphing,decorations.pathreplacing,backgrounds,shadows,babel}
\pgfplotsset{compat=1.18}
\usepgfplotslibrary{fillbetween}
\usepackage[np]{numprint}
\usepackage{listings,makecell,pas-tableur,twoopt}
\let\barre\undefined
\usepackage[pstricks]{tablvar}
\definecolor{darkgreen}{RGB}{0,100,0}
\providecommand{\red}{\color{red}}
\providecommand{\blue}{\color{blue}}
\providecommand{\green}{\color{green}}
\providecommand{\black}{\color{black}}
\providecommand{\white}{\color{white}}
\providecommand{\gray}{\color{gray}}
\providecommand{\cyan}{\color{cyan}}
\providecommand{\magenta}{\color{magenta}}
\providecommand{\yellow}{\color{yellow}}
\providecommand{\violet}{\color{violet}}
\providecommand{\orange}{\color{orange}}
\providecommand{\euro}{\texteuro}
\providecommand{\EUR}{\texteuro}
\providecommand{\ieme}{\textsuperscript{e}}
\providecommand{\degres}{\ensuremath{^\circ}}
\providecommand{\degree}{\ensuremath{^\circ}}
\providecommand{\vv}[1]{\overrightarrow{#1}}
\providecommand{\Coord}[1]{#1}
\providecommand{\decosix}{}
\providecommand{\starredbullet}{\ensuremath{\bullet}}
\pagestyle{empty}
\setlength{\parindent}{0pt}
\usepackage{adjustbox}\DeclareUnicodeCharacter{2236}{:}\DeclareUnicodeCharacter{2190}{\ensuremath{\leftarrow}}
\providecommand{\R}{\mathbb{R}}
\providecommand{\N}{\mathbb{N}}
\providecommand{\Z}{\mathbb{Z}}
\providecommand{\Q}{\mathbb{Q}}
\providecommand{\D}{\mathbb{D}}
\providecommand{\C}{\mathbb{C}}
\providecommand{\e}{{\mathrm{e}}}
\providecommand{\vect}[1]{\overrightarrow{#1}}
\providecommand{\vectt}[1]{\overrightarrow{\mathrm{#1}}}
\providecommand{\vvt}[1]{\overrightarrow{\mathrm{#1}}}
\providecommand{\np}[2][]{{#2~#1}}
\providecommand{\vv}[1]{{\overrightarrow{#1}}}
\renewcommand{\d}{{\mathrm{d}}}
\providecommand{\degres}{{^\circ}}
\providecommand{\degre}{{^\circ}}
\providecommand{\blue}{\color{blue}}
\providecommand{\red}{\color{red}}
\providecommand{\enskip}{\quad}
\providecommand{\llbracket}{\lbrack\!\lbrack}
\providecommand{\rrbracket}{\rbrack\!\rbrack}
\providecommand{\euro}{{\text{€}}}
\providecommand{\no}{{\text{n°}}}
\providecommand{\interval}[2]{{[#1;#2]}}
\providecommand{\intervFO}[2]{{[#1;#2[}}
\providecommand{\dotfill}{\ldots}

\providecommand{\ds}{\displaystyle}
\providecommand{\textding}[1]{#1}
\providecommand{\Oij}{\(\mathrm{O};\vec{\imath},\vec{\jmath}\)}
\providecommand{\Oijk}{\(\mathrm{O};\vec{\imath},\vec{\jmath},\vec{k}\)}
\providecommand{\Ouv}{\(\mathrm{O};\vec{u},\vec{v}\)}
\providecommand{\intervFF}[2]{\ensuremath{[#1;#2]}}
\renewcommand{\euro}{\ifmmode\text{\texteuro}\else\texteuro\fi}
\DeclareUnicodeCharacter{00B0}{\ifmmode^{\circ}\else\textdegree\fi}
\DeclareUnicodeCharacter{20AC}{\ifmmode\text{\texteuro}\else\texteuro\fi}
% Énoncés : Éducation nationale. Transcription et corrigés : APMEP, Denis Vergès et contributeurs. Découpage et classement : CoopMaths, Sébastien Lozano, Éric Elter et contributeurs. Adaptation de la présentation : Maths Explorer.
% CC BY-SA : https://coopmaths.fr/www/about/
% https://coopmaths.fr/alea/static/e3c/2020/tex/e3c_2020_00_sujet51_1.tex
\begin{document}
\noindent\textit{APMEP / CoopMaths. CC BY-SA. Présentation adaptée par Maths Explorer.}\par\medskip
\medskip


\emph{Pour chacune des questions, une seule des quatre réponses proposées est correcte.}

\emph{Les questions sont indépendantes.}

\emph{Pour chaque question, indiquer le numéro de la question et recopier sur la copie la lettre correspondant à la réponse choisie.}

\emph{Aucune justification n'est demandée, mais il peut être nécessaire d'effectuer des recherches au brouillon pour aider à déterminer votre réponse.}

\emph{Chaque réponse correcte rapporte un point. Une réponse incorrecte ou une question sans réponse n'apporte ni ne retire aucun point.}

\medskip

\textbf{Question 1}

\medskip 

Dans un repère orthonormé, le cercle de centre A$(2~;~-1)$ et de rayon 4 a comme équation :

\medskip

\begin{tabularx}{\linewidth}{*{4}{X}}
\textbf{a.~~} $(x+2)^2+(y-1)^2=16$ &\textbf{b.~~} $(x-2)^2+(y+1)^2=4$&\textbf{c.~~}$(x-2)^2+(y+1)^2=16$& \textbf{d.~~} $(x+2)^2+(y-1)^2=4$.
\end{tabularx}

\medskip
\textbf{Question 2}
\medskip 

Soit la droite (d) d'équation cartésienne $2x-y+1 = 0$.
Sachant que la droite $(d_1)$ est perpendiculaire à la droite (d), une équation de $(d_1)$
 peut être :

\begin{tabularx}{\linewidth}{*{4}{X}}
\textbf{a.~~} $x-2y+2 = 0$ &\textbf{b.~~} $x+2y-1 = 0$&\textbf{c.~~}$-2x+y-1 = 0$& \textbf{d.~~} $x-y+2 = 0$.
\end{tabularx}

\medskip
\textbf{Question 3}
\medskip 
L'expression de $\sin(\pi- x)+\cos \left(x+ \frac{\pi}{2}\right)$ est égale à :


\begin{tabularx}{\linewidth}{*{4}{X}}
\textbf{a.~~} $-2\sin(x)$ &\textbf{b.~~} $0$&\textbf{c.~~}$2\sin(x)$& \textbf{d.~~} $\cos(x)-\sin (x)$.
\end{tabularx}

\medskip
\textbf{Question 4}
\medskip 

On considère la fonction polynôme du second degré $f$ définie sur $\R$ par $f(x)= -3x^2+x-5$

Le tableau de variations de cette fonction est :

\begin{tabularx}{\linewidth}{cXcX}
\textbf{a.~~}&
\begin{pspicture}(6.75,2.8)
\psframe(6.65,2.7)
\psline(0,1.7)(6.65,1.7)  \psline(1.2,0)(1.2,2.7)
\uput[u](0.7,2){$x$} \uput[u](1.8,2){$-\infty$}\uput[u](3.65,1.73){$-\dfrac{1}{6}$}\uput[u](6.3,2){$+\infty$}
\uput[d](0.7,1.2){$f$}
\psline{->}(2,0.2)(3.4,1.25)\psline{->}(4.65,1.25)(5.8,0.2)
\end{pspicture}
&\textbf{b.~~}&
\begin{pspicture}(6.75,2.8)
\psframe(6.65,2.7)
\psline(0,1.7)(6.65,1.7)  \psline(1.2,0)(1.2,2.7)
\uput[u](0.7,2){$x$} \uput[u](1.8,2){$-\infty$}\uput[u](3.65,1.73){$\dfrac{1}{3}$}\uput[u](6.3,2){$+\infty$}
\uput[d](0.7,1.2){$f$}
\psline{->}(2,0.2)(3.4,1.25)\psline{->}(4.65,1.25)(5.8,0.2)
\end{pspicture} \\
\textbf{c.~~}&
\begin{pspicture}(6.75,2.8)
\psframe(6.65,2.7)
\psline(0,1.7)(6.65,1.7)  \psline(1.2,0)(1.2,2.7)
\uput[u](0.7,2){$x$} \uput[u](1.8,2){$-\infty$}\uput[u](3.65,1.73){$-\dfrac{1}{3}$}\uput[u](6.3,2){$+\infty$}
\uput[d](0.7,1.2){$f$}
\psline{->}(2,1.25)(3.4,0.2)\psline{->}(4.65,0.2)(5.8,1.25)
\end{pspicture}&
\textbf{d.~~}&
\begin{pspicture}(6.75,2.8)
\psframe(6.65,2.7)
\psline(0,1.7)(6.65,1.7)  \psline(1.2,0)(1.2,2.7)
\uput[u](0.7,2){$x$} \uput[u](1.8,2){$-\infty$}\uput[u](3.65,1.73){$\dfrac{1}{6}$}\uput[u](6.3,2){$+\infty$}
\uput[d](0.7,1.2){$f$}
\psline{->}(2,0.2)(3.4,1.25)\psline{->}(4.65,1.25)(5.8,0.2)
\end{pspicture}.
\end{tabularx}

\medskip

\textbf{Question 5}

\medskip

À un jeu, la variable aléatoire donnant le gain algébrique $G$ suit la loi de probabilité suivante (en euros) :

\begin{center}
\begin{tabular}[]{|m{2cm}|*{4}{c|}}
\hline
Valeurs de $G$&$-25$&$-3$&$x$&$100$\\\hline
Probabilité&$\dfrac{1}{3}$&$\dfrac{1}{6}$&0,3&0,2\rule[-3mm]{0mm}{9mm}\\\hline
\end{tabular}
\end{center}

Sachant que l'espérance de $G$ est égale à $\dfrac{38}{3}$, la valeur de $x$ est :

\begin{tabularx}{\linewidth}{*{4}{X}}
\textbf{a.~~} $0$ &\textbf{b.~~} $5$&\textbf{c.~~}$20$& \textbf{d.~~} $25$.
\end{tabularx}

\vspace{0,5cm}

\end{document}
